\scanpage{5}
% Two short horizontal strokes above the first line.
\margin{\mbox{-- --}}

donc $E'(M\;N)$ est un complexe dont les invariants $\underline{H}^{i}$ sont ceux de $E(M,N)$ en degré $i\unsure{+|\ne}2$, et qui en degré $2$ donne lieu \unsure{:|à} une suite exacte
\begingroup
\makeatletter\tagsleft@true\makeatother
\[
\tag*{(\unsure{S|*})}
0\to\underline{\mathrm{Ext}}^{2}(M,N)\to
\overbrace{\underline{H}^{2}\!\left(E'(M,N)\right)}^{P(M,N)}
\xrightarrow{\sigma}\underline{\mathrm{Hom}}(M,{}_{2}N)\to \unsure{0}
\]
\endgroup

Heuristiquement, $E'(M,N)$ est le complexe qui exprime le ``2-champ de Picard\deleted{\unsure{''|*}} strict'' formé des 1-champs de Picard (pas \unsure{nécessairement} stricts) \inserted{``épinglés'' par $M,N$} sur des objets variables de $X$, en admettant que ta théorie pour les 1-champs de Picard \inserted{stricts} s'étend aux 2-champs de Picard stricts (ce qui pour \unsure{,|.} moi ne fait guère de doute) ; de même $E(M,N)$ correspond aux champs de Picard \uline{stricts} épinglés par $M,N$.

La suite exacte $(*)$ se construit en \unsure{tous|tout} cas canoniquement ``à la main'', où \unsure{le|\deleted{le}} terme médian est le faisceau des \unsure{classes} \inserted{à} \deleted{\unsure{d'|à}}``équivalence'' près des

% Two curved impressions run along the lower left edge of the copied sheet.
% archive sheet number: 62
